\documentclass[12pt,reqno]{amsart}
\usepackage{math-review}
\providecommand{\Ric}{\operatorname{Ric}}
\providecommand{\Hess}{\operatorname{Hess}}
\providecommand{\tr}{\operatorname{tr}}
\providecommand{\Sym}{\operatorname{Sym}}

\title{Completing a block quadratic form}
\date{}
\begin{document}
\maketitle

The question is whether a coupled quadratic form is strictly positive.
We eliminate its mixed term, read off the residual form, and then use the
same calculation to minimize over one block of variables.

\section{A local positivity problem}
The reader needs only matrix multiplication, quadratic-form positivity
and an invertible change of variables. The Schur complement will emerge
from the calculation rather than serve as an assumed rule.

\begin{theorem}\label{thm:schur}
Let $p,q\geq1$, let $A\in\mathbb R^{p\times p}$ and $C\in\mathbb R^{q\times q}$
be symmetric, and let $B\in\mathbb R^{p\times q}$. Assume $A>0$, where
$H>0$ means $v^THv>0$ for every nonzero real vector $v$ of the appropriate size.
Set
\[
 M=\begin{pmatrix}A&B\\B^T&C\end{pmatrix},\qquad S=C-B^TA^{-1}B.
\]
Then $M>0$ if and only if $S>0$. Moreover, for every fixed $y\in\mathbb R^q$,
\[
 \min_{x\in\mathbb R^p}(x,y)^TM(x,y)=y^TSy,
 \qquad x_{\min}=-A^{-1}By.
\]
The minimizer is unique. This is a classical quadratic-form identity and criterion.
\end{theorem}

\subsection*{The point: cancel the mixed term}
For fixed blocks in the theorem, put
\[
 Q(x,y)=x^TAx+2x^TBy+y^TCy.
\]
The mixed term will occur in a square if the new variable $z$ satisfies
$Az=Ax+By$. Since $A>0$ implies that $A$ is invertible, choose
\begin{equation}\label{eq:change}
 z=x+A^{-1}By,\qquad x=z-A^{-1}By.
\end{equation}
Symmetry gives $(A^{-1})^T=A^{-1}$. Expanding without discarding any term,
\[
 z^TAz=x^TAx+2x^TBy+y^TB^TA^{-1}By.
\]
Subtracting the last term from the $y$-part produces
\begin{equation}\label{eq:square}
 Q(x,y)=z^TAz+y^T\underbrace{(C-B^TA^{-1}B)}_{S}y.
\end{equation}
The matrix $S$ is called the Schur complement of $A$. Its role is now
visible: it is exactly the residual quadratic form after cancellation
of the coupling; see \cite[Appendix A.5.5]{S-Boyd}.

\subsection*{The connection: both implications}
The two formulas in \eqref{eq:change} show that $(x,y)\mapsto(z,y)$ is a
bijection preserving the nonzero-pair condition. If $S>0$, each nonzero
pair has a nonzero component, so \eqref{eq:square} is strictly positive.

Conversely, suppose $M>0$. For each $y\ne0$, take $x=-A^{-1}By$.
Then $(x,y)\ne0$ while $z=0$, and
\[
 0<Q(-A^{-1}By,y)=y^TSy.
\]
This proves $S>0$ because the test applies to every nonzero $y$.
The argument uses positivity of a quadratic form, without assuming
that the entire block matrix is invertible in advance.

\section{The minimum and the limit of the argument}
For a fixed $y$, the first term of \eqref{eq:square} is nonnegative and
vanishes only when $z=0$. The inverse change therefore gives the unique
minimum in Theorem~\ref{thm:schur}.

For the scalar blocks $A=2$, $B=1$, $C=1$, the calculation is concrete:
\[
 Q(x,y)=2x^2+2xy+y^2
       =2\left(x+\frac y2\right)^2+\frac{y^2}{2}.
\]
Thus $S=1/2>0$, and the minimizing choice is $x=-y/2$. Both positivity
and the minimum value can be checked in the same displayed expression.

If $A$ is symmetric and merely invertible, the square identity still
holds. The inequality needed to minimize need not hold. For example,
\[
 A=-1,\quad B=0,\quad C=1
 \quad\Longrightarrow\quad S=1,\quad Q(x,y)=-x^2+y^2.
\]
Here $S$ is positive but $M$ is indefinite, and $Q(\cdot,y)$ is unbounded
below. If $A$ is singular, $A^{-1}$ is unavailable; generalized criteria
require additional work and are not asserted here.

\bibliographystyle{../manuscript/localrefs}
\bibliography{../generated/references}
\end{document}
